% GENERATED FILE. Edit canonical lessons, then run npm run generate:books.
% canonical-source-hash: fnv1a64:8b5ebda86d07ee3f
% canonical-lessons: SA-C96-rajamargah, SA-C96-catuspathah, SA-C96-konah, SA-C96-pustakalayah, SA-C96-visvavidyalayah

\chapter{Main road, Crossroads, Corner, Library, University}
\label{ch:sa-main-road-crossroads-corner-library-university}
\hlchaptermodality{96}

\begin{chapteropening}
\textbf{By the end of this chapter:} \emph{I can say a main road, a crossroads, a corner, a library and a university.}

Complete the last lesson of chapter 96: I can say a main road, a crossroads, a corner, a library and a university.
\end{chapteropening}

\section[\texorpdfstring{\sk{राजमार्गः} (rājamārgaḥ)}{राजमार्गः (rājamārgaḥ)}]{\sk{राजमार्गः} (rājamārgaḥ) — a main road}

\label{lesson:SA-C96-rajamargah}



\begin{quote}
Before the new one: say the Sanskrit for elsewhere, then the Sanskrit for a street.
\end{quote}

\subsection*{You'll want to know: \sk{राजमार्गः}}
\textbf{\sk{राजमार्गः}} — \emph{rājamārgaḥ} — \textquotedblleft{}a main road\textquotedblright{}.

\begin{cousinweb}
\emph{rājamārga} is a compound of \emph{rājan} \textquotedblleft{}king\textquotedblright{} and \emph{mārga} \textquotedblleft{}road\textquotedblright{}: \textquotedblleft{}the king's highway\textquotedblright{}, a main street wide enough for horses and elephants (MW). \emph{mārga} comes from \emph{mṛga} \textquotedblleft{}deer, game\textquotedblright{} (MW), and one of its senses is the track of a wild animal.
\end{cousinweb}

\begin{grammarlens}[title={What you've built}]
One more word for this chapter.
\end{grammarlens}

\subsection*{Your turn}
\begin{itemize}
\raggedright
  \item \emph{Say it:} \emph{rājamārgaḥ}
  \item \emph{Say it:} \emph{rājamārgaḥ}, once more
  \item \emph{Say it:} say \emph{rājamārgaḥ}
  \item \emph{Recall:} say the Sanskrit for elsewhere, then the Sanskrit for a street, then say \emph{rājamārgaḥ} again
\end{itemize}

\begin{tcolorbox}[breakable,colback=teal!4,colframe=teal!35!black,title={Before you move on}]
What does \sk{राजमार्गः} mean? (\textquotedblleft{}A main road\textquotedblright{}.) Say it once more.
\end{tcolorbox}
\section[\texorpdfstring{\sk{चतुष्पथः} (catuṣpathaḥ)}{चतुष्पथः (catuṣpathaḥ)}]{\sk{चतुष्पथः} (catuṣpathaḥ) — a crossroads}

\label{lesson:SA-C96-catuspathah}



\begin{quote}
Before the new one: say the Sanskrit for a street, then the Sanskrit for a main road.
\end{quote}

\subsection*{You'll want to know: \sk{चतुष्पथः}}
\textbf{\sk{चतुष्पथः}} — \emph{catuṣpathaḥ} — \textquotedblleft{}a crossroads\textquotedblright{}.

\begin{cousinweb}
\emph{catuṣpatha} is a compound of \emph{catur} \textquotedblleft{}four\textquotedblright{} and \emph{patha} \textquotedblleft{}road\textquotedblright{}: \textquotedblleft{}where four roads meet\textquotedblright{}, already in the Brāhmaṇas (MW). \emph{catur} is the inherited number behind Latin \emph{quattuor} (MW) and English \emph{four}. Hindi \emph{caupath} \textquotedblleft{}crossroads\textquotedblright{} comes from the same pair (CDIAL 4642).
\end{cousinweb}

\begin{grammarlens}[title={What you've built}]
One more word for this chapter.
\end{grammarlens}

\subsection*{Your turn}
\begin{itemize}
\raggedright
  \item \emph{Say it:} \emph{catuṣpathaḥ}
  \item \emph{Say it:} \emph{catuṣpathaḥ}, once more
  \item \emph{Say it:} say \emph{catuṣpathaḥ}
  \item \emph{Recall:} say the Sanskrit for a street, then the Sanskrit for a main road, then say \emph{catuṣpathaḥ} again
\end{itemize}

\begin{tcolorbox}[breakable,colback=teal!4,colframe=teal!35!black,title={Before you move on}]
What does \sk{चतुष्पथः} mean? (\textquotedblleft{}A crossroads\textquotedblright{}.) Say it once more.
\end{tcolorbox}
\section[\texorpdfstring{\sk{कोणः} (koṇaḥ)}{कोणः (koṇaḥ)}]{\sk{कोणः} (koṇaḥ) — a corner}

\label{lesson:SA-C96-konah}



\begin{quote}
Before the new one: say the Sanskrit for a main road, then the Sanskrit for a crossroads.
\end{quote}

\subsection*{You'll want to know: \sk{कोणः}}
\textbf{\sk{कोणः}} — \emph{koṇaḥ} — \textquotedblleft{}a corner\textquotedblright{}.

\begin{cousinweb}
\emph{koṇa} \textquotedblleft{}corner\textquotedblright{} is often taken to be borrowed from Dravidian: CDIAL 3504 follows Burrow in deriving it from a Dravidian word, and DEDR 2209 sets it beside a Dravidian family of words for corners. It lives on as Hindi \emph{konā} \textquotedblleft{}corner\textquotedblright{} (CDIAL 3504).
\end{cousinweb}

\begin{grammarlens}[title={What you've built}]
One more word for this chapter.
\end{grammarlens}

\subsection*{Your turn}
\begin{itemize}
\raggedright
  \item \emph{Say it:} \emph{koṇaḥ}
  \item \emph{Say it:} \emph{koṇaḥ}, once more
  \item \emph{Say it:} say \emph{koṇaḥ}
  \item \emph{Recall:} say the Sanskrit for a main road, then the Sanskrit for a crossroads, then say \emph{koṇaḥ} again
\end{itemize}

\begin{tcolorbox}[breakable,colback=teal!4,colframe=teal!35!black,title={Before you move on}]
What does \sk{कोणः} mean? (\textquotedblleft{}A corner\textquotedblright{}.) Say it once more.
\end{tcolorbox}
\section[\texorpdfstring{\sk{पुस्तकालयः} (pustakālayaḥ)}{पुस्तकालयः (pustakālayaḥ)}]{\sk{पुस्तकालयः} (pustakālayaḥ) — a library}

\label{lesson:SA-C96-pustakalayah}



\begin{quote}
Before the new one: say the Sanskrit for a crossroads, then the Sanskrit for a corner.
\end{quote}

\subsection*{You'll want to know: \sk{पुस्तकालयः}}
\textbf{\sk{पुस्तकालयः}} — \emph{pustakālayaḥ} — \textquotedblleft{}a library\textquotedblright{}.

\begin{cousinweb}
\emph{pustakālaya} is a compound of \emph{pustaka} \textquotedblleft{}book\textquotedblright{} and \emph{ālaya} \textquotedblleft{}dwelling, house\textquotedblright{} (MW); the pairing for a library is modern. \emph{pustaka} was borrowed from Iranian: CDIAL 8413 compares Sogdian \emph{pwstk} \textquotedblleft{}book\textquotedblright{} and Persian \emph{pōst} \textquotedblleft{}skin\textquotedblright{}.
\end{cousinweb}

\begin{grammarlens}[title={What you've built}]
One more word for this chapter.
\end{grammarlens}

\subsection*{Your turn}
\begin{itemize}
\raggedright
  \item \emph{Say it:} \emph{pustakālayaḥ}
  \item \emph{Say it:} \emph{pustakālayaḥ}, once more
  \item \emph{Say it:} say \emph{pustakālayaḥ}
  \item \emph{Recall:} say the Sanskrit for a crossroads, then the Sanskrit for a corner, then say \emph{pustakālayaḥ} again
\end{itemize}

\begin{tcolorbox}[breakable,colback=teal!4,colframe=teal!35!black,title={Before you move on}]
What does \sk{पुस्तकालयः} mean? (\textquotedblleft{}A library\textquotedblright{}.) Say it once more.
\end{tcolorbox}
\section[\texorpdfstring{\sk{विश्वविद्यालयः} (viśvavidyālayaḥ)}{विश्वविद्यालयः (viśvavidyālayaḥ)}]{\sk{विश्वविद्यालयः} (viśvavidyālayaḥ) — a university}

\label{lesson:SA-C96-visvavidyalayah}



\begin{quote}
Before the new one: say the Sanskrit for a corner, then the Sanskrit for a library.
\end{quote}

\subsection*{You'll want to know: \sk{विश्वविद्यालयः}}
\textbf{\sk{विश्वविद्यालयः}} — \emph{viśvavidyālayaḥ} — \textquotedblleft{}a university\textquotedblright{}.

\begin{cousinweb}
\emph{viśvavidyālaya} is a compound of \emph{viśva} \textquotedblleft{}all, whole\textquotedblright{} and \emph{vidyālaya} \textquotedblleft{}seat of learning\textquotedblright{} (MW), itself \emph{vidyā} \textquotedblleft{}knowledge\textquotedblright{} + \emph{ālaya} \textquotedblleft{}dwelling\textquotedblright{}; using it for a university is modern. \emph{vidyā} comes from the root √vid \textquotedblleft{}know\textquotedblright{}, which MW sets beside Latin \emph{videre} \textquotedblleft{}see\textquotedblright{} (English \emph{video}) and English \emph{wot} \textquotedblleft{}know\textquotedblright{}.
\end{cousinweb}

\begin{grammarlens}[title={What you've built}]
One more word for this chapter.
\end{grammarlens}

\subsection*{Your turn}
\begin{itemize}
\raggedright
  \item \emph{Say it:} \emph{viśvavidyālayaḥ}
  \item \emph{Say it:} \emph{viśvavidyālayaḥ}, once more
  \item \emph{Say it:} say \emph{viśvavidyālayaḥ}
  \item \emph{Recall:} say the Sanskrit for a corner, then the Sanskrit for a library, then say \emph{viśvavidyālayaḥ} again
\end{itemize}

\begin{tcolorbox}[breakable,colback=teal!4,colframe=teal!35!black,title={Before you move on}]
What does \sk{विश्वविद्यालयः} mean? (\textquotedblleft{}A university\textquotedblright{}.) Say it once more.
\end{tcolorbox}
